\title{FlashAttention-3: Fast and Accurate Attention with Asynchrony and Low-precision}

\begin{abstract}
Attention serves as a foundational layer within the widely-adopted Transformer architecture, and remains the primary performance constraint for scaling large language models and supporting long-context sequences. \textsc{FlashAttention} previously introduced a method to accelerate attention computation on GPUs by reducing the frequency of memory accesses. Despite this, the method does not fully leverage the enhanced functionalities in state-of-the-art hardware, with \textsc{FlashAttention-2} utilizing just 35\% of the H100 GPU's processing capability. To address these limitations, we introduce three strategies tailored for Hopper GPUs: (1) employing the asynchronous capabilities of Tensor Cores and TMA to overlap both computation and data transfers through warp-specialization, (2) interlacing block-level matrix multiplication with softmax computations, and (3) integrating block quantization with incoherent operations to exploit hardware-enabled FP8 low-precision arithmetic. Our proposed approach, \textsc{FlashAttention-3}, delivers a 1.5-2.0$\times$ throughput increase on H100 GPUs, achieving up to 740 TFLOPs/s with FP16 (amounting to 75\% utilization) and approaching 1.2 PFLOPs/s with FP8 precision. Additionally, our evaluation confirms that FP8 \textsc{FlashAttention-3} results in 2.6$\times$ reduction in numerical error compared to standard FP8 attention implementations.
\end{abstract}

\section{Introduction}
\label{sec:intro}

Within the Transformer architecture~\citep{vaswani2017attention}, the attention mechanism is the central source of computational expense, primarily because the process of calculating self-attention scores between queries and keys requires computation that increases quadratically with respect to the sequence length. Enabling attention mechanisms to efficiently handle substantially longer sequences could open up new possibilities, such as advanced modeling and reasoning across multiple extensive documents~\citep{guo2021longt5,shaham2022scrolls,peng2023yarn}, or navigating through large-scale codebases~\citep{roziere2023code, li2023starcoder}. Expanded capabilities may also include support for new modalities like high-resolution imagery~\citep{chen2022scaling}, audio~\citep{gulati2020conformer}, and video content~\citep{ho2022video}. Moreover, it may enable novel applications, such as managing user interactions over lengthy histories~\citep{sun2019bert4rec} or agent workflows with extended planning horizons~\citep{yao2022react}. This demand for scalability has spurred intensive research into accelerating attention for longer contexts, leveraging both approximation techniques~\citep{katharopoulos2020transformers,choromanski2020rethinking,
tay2020efficient}, improved software implementations (\citep{rabe2021self,dao2022flashattention,kwon2023efficient}), and even alternative model architectures~\citep{peng2023rwkv,sun2023retentive,gu2023mamba}.

This study extends the advancements presented by \citet{dao2022flashattention}, who proposed exact-attention algorithms informed by GPU execution paradigms and underlying hardware traits. Their approach in \citep{dao2022flashattention}, known as \textsc{FlashAttention}, relies on a distinct tiling scheme that enables efficient parallelization of attention operations by consolidating these steps into a single GPU kernel. This fusion effectively removes extraneous accesses to the slower global memory.

Building on this, \citet{dao2023flashattention2} introduced \textsc{FlashAttention-2}, redesigning the algorithm to include sequence length dimension parallelism and block-wise processing of key and value matrices within the forward pass. These enhancements aim to elevate GPU resource utilization by better distributing computational tasks. Nevertheless, our analysis reveals that \textsc{FlashAttention-2} is still suboptimal in terms of hardware usage on modern GPUs—achieving only approximately 35\% occupancy compared to the 80–90\% demonstrated by highly-tuned matrix-multiplication (GEMM) routines on the Hopper H100 GPU. This inefficiency is partly due to practical aspects of implementation, such as relying on Ampere-specific instructions rather than leveraging the specialized Hopper instructions within Tensor Core operations. Recent work, including ThunkerKitten~\citep{spector2024thunder} and cuDNN 9~\citep{cudnn9}, has demonstrated that employing Hopper-native instructions and tile-based methods results in both more efficient attention calculations and a more streamlined programming process.

At its core, the algorithm utilized by \textsc{FlashAttention-2} operates within a straightforward synchronous framework, lacking explicit integration of asynchronous execution or low-precision arithmetic in its approach. Asynchrony emerges from specialized hardware designed to optimize critical machine learning operations: dedicated units are assigned to matrix multiplications (Tensor Cores) or memory access (Tensor Memory Accelerator -- TMA), independently of standard CUDA cores responsible for logical operations, integer math, and floating-point computations. The adoption of low-precision formats, exemplified by FP8 in Hopper and FP4 in Blackwell—with earlier precedents such as FP16 (Pascal, 2017) and BF16 (Ampere, 2020)—has consistently demonstrated increased computational throughput, providing two or four times higher performance without expanding power consumption or die size. \cref{subsec:hardware} surveys the advances introduced in Hopper along these axes.

A central engineering task is to rearchitect \textsc{FlashAttention-2} so it fully leverages these hardware innovations: exploiting asynchrony requires effectively interleaving matmul and softmax computations despite their sequential dependency, while utilizing low-precision demands precise strategies to control quantization error, particularly given the occurrence of outlier features in large language models~\citep{dettmers2208llm, sun2024massive}.

In pursuit of this objective, we introduce \textsc{FlashAttention-3}, which integrates and advances three innovative techniques to enhance performance on modern GPU architectures:\footnote{Our discussion centers around NVIDIA's Hopper architecture, but the algorithm remains effective across GPU architectures featuring adequate asynchronous execution and support for low-precision computations.}

\begin{enumerate}

\item \textbf{Producer-Consumer asynchrony:} We introduce a warp-specialized form of software pipelining that leverages the asynchronous operation of data transfer and Tensor Cores. By dividing data producers and consumers across distinct warps, our approach enhances the algorithm's capacity to obscure memory access and instruction issuance delays.
\item \textbf{Hiding softmax under asynchronous block-wise GEMMs:} The softmax procedure, which consists of less computationally intensive routines such as floating point multiply-add and exponential functions, is executed concurrently with asynchronous WGMMA instructions for GEMM. To facilitate this, we modify the \textsc{FlashAttention-2} methodology, eliminating select sequential interdependencies between softmax and GEMM computations. As an illustration, in our algorithm's 2-stage variant, softmax is performed for one score block while WGMMA simultaneously handles the asynchronous calculation of the subsequent block.
\item \textbf{Hardware-accelerated low-precision GEMM:} We adjust the forward propagation process to utilize FP8 Tensor Cores for GEMM, effectively achieving nearly twice the throughput in TFLOPs/s. Accommodating this requires addressing the layout differences dictated by WGMMA regarding FP32 accumulators and FP8 operand matrices stored in memory. By employing block quantization and incoherent processing, we minimize the degradation in accuracy associated with the shift to FP8 computation.
\end{enumerate}

In order to empirically substantiate our approach, we systematically evaluate \textsc{FlashAttention-3} on the H100 SXM5 GPU across multiple parameter configurations and report the following: (1) FP16 attains a speed increase of 1.5-2.0$\times$ compared to \textsc{FlashAttention-2} during the forward pass (peaking at 740 TFLOPs/s) and delivers a 1.5-1.75$\times$ acceleration in the backward pass; (2) FP8 registers performance near 1.2 PFLOPs/s; and (3) for extended sequence lengths, FP16 demonstrates superior throughput while FP8 remains competitive\footnote{Specifically, \textsc{FlashAttention-3} FP8 leads for head dimension 64; for head dimensions 128 and 256, FP8 matches performance absent causal masking, but lags when causal masking is enabled.} with NVIDIA cuDNN's leading attention implementation. Additionally, we confirm that FP16 \textsc{FlashAttention-3} maintains numerical precision equivalent to \textsc{FlashAttention-2}, surpassing the conventional attention computation where intermediate operations (such as softmax rescaling) are carried out in FP32. Furthermore, for cases with outlier features, FP8 \textsc{FlashAttention-3} using block quantization and incoherent processing achieves 2.6$\times$ greater accuracy relative to standard attention employing per-tensor quantization.

We have released \textsc{FlashAttention-3} as open-source software under a liberal license\footnote{\textsc{FlashAttention-3} is accessible at \texttt{https://github.com/Dao-AILab/flash-attention}}. Our aim is to incorporate it into the PyTorch and Hugging Face ecosystems to maximize its utility for researchers and developers.

\section{Background: Multi-Head Attention and GPU Characteristics}
\label{sec:background}

\subsection{Multi-Head Attention}
\label{subsec:multi_head_attn}

Let $\mathbf{Q}, \mathbf{K}, \mathbf{V} \in \mathbb{R}^{N \times d}$ be the query, key and value input sequences associated to a single head, where $N$ is the sequence length and $d$ is the head dimension. Then the attention output $\mathbf{O}$ is computed as:

\begin{equation*}
  \mathbf{S} = \alpha \mathbf{Q} \mathbf{K}^\top \in \mathbb{R}^{N \times N}, \quad \mathbf{P} = \mathrm{softmax}(\mathbf{S}) \in \mathbb{R}^{N \times N}, \quad \mathbf{O} = \mathbf{P}\mathbf{V} \in \mathbb{R}^{N \times d},
\end{equation*}

Here, the $\mathrm{softmax}$ function is computed for each row, and it is common to use $\alpha = 1/\sqrt{d}$ as the scale parameter. To address potential numerical issues arising from the exponential operation, we typically subtract $\mathrm{rowmax}(\mathbf{S})$ from $\mathbf{S}$. In the context of multi-head attention (MHA), individual heads maintain distinct query, key, and value projection matrices. The operations are performed in parallel over various heads and batch instances to generate the complete output tensor.

Now let $\phi$ be a scalar loss function and let $\mathbf{d}(-) = \partial \phi / \partial (-)$ be notation for the gradient. Given the output gradient $\mathbf{dO} \in \mathbb{R}^{N \times d}$, we compute $\mathbf{dQ}$, $\mathbf{dK}$, and $\mathbf{dV}$ according to the chain rule as follows:

\begin{align*}
  \mathbf{dV} &= \mathbf{P}^\top \mathbf{dO} \in \mathbb{R}^{N \times d} \\
  \mathbf{dP} &= \mathbf{dO} \mathbf{V}^\top \in \mathbb{R}^{N \times N} \\
  \mathbf{dS} &= \mathrm{dsoftmax} (\mathbf{dP}) \in \mathbb{R}^{N \times N} \\
  \mathbf{dQ} &= \alpha \mathbf{dS} \mathbf{K} \in \mathbb{R}^{N \times d} \\
  \mathbf{dK} &= \alpha \mathbf{dS}^\top \mathbf{Q} \in \mathbb{R}^{N \times d},
\end{align*}

In this context, $\mathbf{d}s = (\mathrm{diag}(p) - p p^\top)\mathbf{d}p$ holds when $p$ represents the $\mathrm{softmax}(s)$ of some input vector $s$, and the notation $\mathrm{dsoftmax}(\mathbf{dP})$ indicates application of this operation to each row. This procedure also allows for efficient parallelization across both heads and batches during the backward propagation step of MHA.

\subsection{GPU hardware characteristics and execution model}
\label{subsec:hardware}

We outline the components of the GPU execution model pertinent to \textsc{FlashAttention-3}, concentrating particularly on the NVIDIA Hopper architecture as a representative example of this model.

\paragraph{Memory hierarchy:} The GPU features a layered structure of memory locations, where storage capacity decreases as bandwidth increases (\cref{tab:gpu-hierarchy})\footnote{\citet{luo2024benchmarking} presents a shared memory bandwidth of 128 bytes per clock per SM; to estimate peak, this is multiplied by 132 SMs and a boost frequency of 1830 MHz.}. The primary global memory (GMEM), also referred to as HBM, resides off-chip and can be accessed by all streaming multiprocessors (SMs). When data is loaded from GMEM, it is automatically stored in an on-chip L2 cache. Proceeding further, every SM integrates a relatively small on-chip cache called shared memory (SMEM), which is managed explicitly by programmers and organized into multiple banks. At the top of the hierarchy, each SM includes its own register file.

\paragraph{Thread hierarchy:} The GPU programming paradigm structures its execution units through a hierarchy of threads. Starting from the most granular level, this hierarchy includes individual threads, followed by warps (each consisting of 32 threads), warpgroups (which are formed by 4 consecutive warps), threadblocks (also known as cooperative thread arrays or CTAs), threadblock clusters (as introduced in Hopper), and grids at the topmost level.

The two hierarchies share a strong connection. Threads belonging to a single CTA are scheduled together onto a common SM, while CTAs within one cluster are jointly assigned to the same GPC. All threads in a CTA can access SMEM directly, but individual threads possess up to 256 private registers (RMEM) each.

\paragraph{Asynchrony and warp-specialization:}

GPUs are designed as throughput-oriented processors that depend on concurrent execution and asynchronous operations to mask latencies related to memory access and computation. In the context of asynchronous data transfers between GMEM and SMEM, Hopper introduces a specialized hardware component known as the Tensor Memory Accelerator (TMA) \cite[\S7.29]{cuda}. Moreover, differing from previous architectures like Ampere, Hopper’s Tensor Core—accessible through the warpgroup-wide WGMMA instruction \cite[\S9.7.14]{ptx}—enables asynchronous execution and can directly fetch operands from shared memory.

The presence of hardware mechanisms for asynchronous operations enables warp-specialized kernels, where different warps within a CTA are assigned dedicated producer or consumer functions, exclusively performing either data transfer or computation tasks. This separation generally enhances the compiler’s capacity to produce highly efficient instruction schedules \citep{warp-specialization-2011}. Furthermore, Hopper introduces the capability to dynamically redistribute register resources among warpgroups using \verb|setmaxnreg| \cite[\S9.7.17.1]{ptx}, allowing warps tasked with MMAs to access a larger portion of RMEM compared to warps assigned to TMA operations—which require only a single thread.

\paragraph{Low-precision number formats:}
\label{sec:low-precision-gpu}
Contemporary GPUs incorporate dedicated hardware modules designed for efficient low-precision arithmetic. For instance, the WGMMA operation enables the FP8 Tensor Cores on Hopper to achieve twice the throughput per SM relative to FP16 or BF16 computations.

To properly utilize FP8 WGMMA, one must be aware of the operand layout requirements. Consider a GEMM operation multiplying $A \times B^{\top}$, where $A$ is of size $M \times K$ and $B$ is $N \times K$. If an operand—either $A$ or $B$—is contiguous along the outer $M$ or $N$ dimension, it is referred to as \emph{mn-major}; if the contiguity is along the inner $K$ dimension, it is known as \emph{k-major}. FP16 WGMMA accepts operands from SMEM in both mn-major and k-major layouts. In contrast, FP8 WGMMA restricts support exclusively to k-major input formats. Additionally, when fusing consecutive GEMM computations (as in attention mechanisms) within a single kernel, mismatches between FP32 accumulator layouts and FP8 operand formats create difficulties for chained FP8 WGMMA invocations.

Within the scope of attention mechanisms, such layout constraints necessitate adjustments to the construction of an FP8 algorithm; these are detailed in \cref{sec:algofp8}.

\subsection{Standard Attention and Flash Attention} In line with the approach of \citet{dao2022flashattention}, we refer to \textbf{standard attention} as a GPU-based attention mechanism that explicitly stores the intermediate matrices $\mathbf{S}$ and $\mathbf{P}$ in HBM. The foundational innovation behind \textsc{FlashAttention} is the use of a localized softmax reduction, which significantly minimizes costly intermediate memory operations by merging the attention calculation into a unified kernel. This local softmax computation is realized by the set of instructions spanning lines \ref{code-ws:softmax_start}-\ref{code-ws:softmax_end} within the consumer mainloop found in \cref{alg:flash3_wgmma_ws_only}, alongside the scaling adjustments applied to segments of $\mathbf{O}$. A straightforward derivation verifying that this operation produces the intended $\mathbf{O}$ output is provided in \cite[\S 2.3.1]{dao2023flashattention2}.

\section{FlashAttention-3: Algorithm}
\label{sec:algo}

This section provides an overview of the \textsc{FlashAttention-3} algorithm. For clarity, we concentrate on the forward pass, leaving the description of the backward pass algorithm to~\cref{sec:algo_ws_bwd}. We start by outlining the procedure for incorporating warp-specialization and a circular SMEM buffer into the foundational framework of \textsc{FlashAttention-2}. Subsequently, we discuss how the asynchronous nature of WGMMA is harnessed to construct a GEMM-softmax pipeline with two overlapping stages. Finally, we detail the adjustments necessary for FP8, encompassing both the requirements for layout compatibility and improvements in numerical accuracy achieved through block quantization and incoherent processing.

\subsection{Producer-Consumer asynchrony through warp-specialization and pingpong scheduling}
\label{sec:algo_ws}

\paragraph{Warp-specialization}
Similar to the approach taken in \textsc{FlashAttention-2}, the forward computation of \textsc{FlashAttention-3} is inherently parallel across the batch dimension, the number of attention heads, as well as the length of the query sequence. For this reason, a high-level description at the CTA granularity is adequate, focusing on the algorithm's processing of a tile $\mathbf{Q}_i$ from the query matrix to produce the associated tile $\mathbf{O}_i$ of the output. For clarity, we initially discuss the warp-specialization strategy employing a circular SMEM buffer, intentionally leaving out the overlap with GEMM-softmax operations. Here, $d$ designates the size of the head, $N$ denotes the total sequence length, and we select a query block size $B_r$ that partitions $\mathbf{Q}$ into $T_r = \lceil \frac{N}{B_r} \rceil$ blocks, written as $\mathbf{Q}_1, .., \mathbf{Q}_{T_r}$.

\begin{algorithm}[H]
    \caption{\small\label{alg:flash3_wgmma_ws_only}\textsc{FlashAttention-3} forward pass \textbf{excluding} intra-consumer overlap -- CTA perspective}
    \begin{algorithmic}[1]
\REQUIRE Inputs: $\mathbf{Q}_i \in \mathbb{R}^{B_r \times d}$ and $\mathbf{K}, \mathbf{V} \in \mathbb{R}^{N \times d}$ resident in HBM, and a designated key block size $B_c$ such that $T_c = \lceil \frac{N}{B_c} \rceil$.
\STATE Establish the pipeline controller for handling synchronization, making use of an $s$-stage circular shared memory buffer.
\IF {participating as a producer warpgroup}
\STATE Free a fixed allocation of registers in preparation.
\STATE Trigger data movement of $\mathbf{Q}_i$ from HBM into shared memory.
\STATE Once transfer is complete, inform consumers that $\mathbf{Q}_i$ has been loaded.
\FOR{$j$ from $0$ up to $T_c-1$}
    \STATE Await confirmation that the $(j\,\%\,s)$ stage in the buffer has been processed.
    \STATE Initiate memory transfers for $\mathbf{K}_j$ and $\mathbf{V}_j$ from HBM into shared memory in buffer stage $(j\,\%\,s)$.
    \STATE Upon loading, signal to consumers that $\mathbf{K}_j$ and $\mathbf{V}_j$ are available.
\ENDFOR
\ELSE \STATE Allocate a register count based on the number of consumer warps.
\STATE Allocate and initialize on-chip buffers: set $\mathbf{O}_i = (0) \in \mathbb{R}^{B_r \times d}$, $\ell_i = 0$, $m_i = -\infty$ in $\mathbb{R}^{B_r}$.
\STATE Block until $\mathbf{Q}_i$ arrives in shared memory.
\FOR{$j$ from $0$ up to $T_c-1$}
\STATE Wait for availability of $\mathbf{K}_j$ in shared memory.
\STATE Execute SS-GEMM, compute $\mathbf{S}_i^{(j)} = \mathbf{Q}_i \mathbf{K}_j^T$. Synchronize as needed. \label{alg:ws_only_gemm1}
\STATE Cache current $m_i$ as $m_i^{\mathrm{old}}$ and update $m_i = \mathrm{max}(m_i^{\mathrm{old}}, \mathrm{rowmax}(\mathbf{S}_i^{(j)}))$. \label{code-ws:softmax_start}
\STATE Evaluate unnormalized probabilities $\widetilde{\mathbf{P}}_i^{(j)} = \mathrm{exp}(\mathbf{S}_i^{(j)} - m_i)$ and update accumulator $\ell_i = \mathrm{exp}(m_i^{\mathrm{old}} - m_i)\ell_i + \mathrm{rowsum}(\widetilde{\mathbf{P}}_i^{(j)})$. \label{code-ws:softmax_end}
\STATE Wait for loading of $\mathbf{V}_j$ in shared memory.
\STATE Compute RS-GEMM: update $\mathbf{O}_i$ via $\mathbf{O}_i = \mathrm{diag}(\mathrm{exp}(m_i^{\mathrm{old}} - m_i))^{-1}\mathbf{O}_i + \widetilde{\mathbf{P}}_i^{(j)}\mathbf{V}_j$. Synchronize afterwards. \label{alg:ws_only_gemm2}
\STATE Mark the $(j\,\%\,s)$ buffer stage as available for the producer.
\ENDFOR
\STATE Normalize: set $\mathbf{O}_i = \mathrm{diag}(\ell_i)^{-1}\mathbf{O}_i$ and compute $L_i = m_i + \log(\ell_i)$.
\STATE Persist $\mathbf{O}_i$ and $L_i$ into HBM, representing the $i$th segments of the output tensors $\mathbf{O}$ and $L$.
\ENDIF
\end{algorithmic}
\end{algorithm}

In our deployment of \cref{alg:flash3_wgmma_ws_only} on Hopper, memory (de)allocation is managed via \verb|setmaxnreg|, while TMA is employed to load $\mathbf{Q}_i$ along with the sets $\{\mathbf{K}_j, \mathbf{V}_j\}_{0 \leq j < T_c}$. The GEMMs executed within the mainloop of the consumer are performed using WGMMA, with the SS or RS prefix denoting whether shared memory or the register file serves as the source for the first operand, respectively. When analyzing the flow of execution for \cref{alg:flash3_wgmma_ws_only}, it is important to recognize that TMA load operations are asynchronous; as a result, the launch of new TMA loads does not block pending loads from finishing. Additionally, within the producer mainloop, no wait instructions will be triggered during the initial $s$ iterations, since these cycles are spent filling the buffer.

\paragraph{Pingpong scheduling} The inherent asynchronicity of both WGMMA and TMA, paired with warp-specialization, allows the possibility for concurrent execution: one warpgroup's softmax computation can proceed while another warpgroup performs its GEMM operation. This is particularly appealing because non-matmul tasks are significantly less performant on current accelerator platforms compared to matmul tasks. For instance, the H100 SXM5 GPU delivers 989 TFLOPS for FP16 matmul operations but only achieves 3.9 TFLOPS for special functions, such as the exponential\footnote{According to the CUDA programming guide, each streaming multiprocessor (SM) can execute 16 special function operations per clock cycle. By multiplying 16 by the total number of SMs (132) and the clock speed (1830 MHz), we arrive at 3.9 TFLOPS for special functions.} required during softmax. Examining the attention forward pass in FP16 with a head dimension of 128, we observe 512 times more FLOPS spent on matmul compared to exponentials; however, the throughput for exponentials is 256 times lower, causing the exponential to occupy approximately 50\% of the total computation time relative to matmul. The discrepancy is further amplified under FP8, where matmul throughput doubles but exponential throughput remains unchanged.

The exponential computation is handled by a dedicated hardware component—the multi-function unit. Optimally, we would want this exponential operation to take place concurrently with the matmul processes executed by the Tensor Cores. To achieve this, we deploy synchronization barriers (\texttt{bar.sync} instructions), orchestrating the GEMM operations such that warpgroup 1 (performing GEMM1 -- $\mathbf{P} \mathbf{V}$ for the current iteration and GEMM0 -- $\mathbf{Q} \mathbf{K}^\top$ for the subsequent iteration) are executed ahead of warpgroup 2's GEMM tasks. This arrangement ensures that the softmax routine for warpgroup 1 coincides with the GEMM execution in warpgroup 2. Afterward, the groups switch, with warpgroup 2 carrying out softmax while warpgroup 1 takes on GEMM calculations—this forms a ``pingpong'' scheduling pattern as depicted in~\cref{fig:pingpong_scheduling}. While in real implementations this pingpong scheduling does not achieve the ideal separation shown in the illustration, empirical results indicate notable performance gains, such as an increase from 570 TFLOPS to between 620 and 640 TFLOPS for FP16 forward passes using a head dimension of 128 and sequence length of 8192.

\paragraph{Attention variants} When utilizing multi-query attention~\citep{shazeer2019fast} and grouped query attention~\citep{ainslie2023gqa}, we adopt the methodology of \textsc{FlashAttention-2}, modifying tensor indexing to prevent unnecessary replication of $\mathbf{K}$ and $\mathbf{V}$ within HBM.

\subsection{Intra-warpgroup overlapping GEMMs and softmax}

Within a single warpgroup, it is possible to execute certain instructions from the softmax concurrently with select instructions from the GEMMs. One method for achieving such overlap is outlined below.

In the attention algorithm, operations within the inner loop (main loop) have sequential dependencies that impede parallelization within a single iteration. For example, (local) softmax (lines \ref{code-ws:softmax_start} to \ref{code-ws:softmax_end}) relies on the output $\mathbf{S}_i^{(j)}$ of the first GEMM, while the second GEMM takes its result $\widetilde{\mathbf{P}}_i^{(j)}$ as an operand. Indeed, the wait statements in lines \ref{alg:ws_only_gemm1} and \ref{alg:ws_only_gemm2} of \cref{alg:flash3_wgmma_ws_only} serialize the execution of softmax and GEMMs. However, we can break these dependencies by pipelining across iterations through additional buffers in registers. Pursuing this idea, we propose the following two-stage\footnote{Note that the number of stages of the overlapping scheme is bounded by, but need not equal, the number $s$ of stages in the circular SMEM buffer.} GEMM-softmax pipelining algorithm:

\begin{algorithm}[H]
    \caption{\small\label{alg:flash3_wgmma}\textsc{FlashAttention-3} consumer warpgroup forward pass}
    \begin{algorithmic}[1]
      \REQUIRE  Matrices $\mathbf{Q}_i \in \mathbb{R}^{B_r \times d}$ and $\mathbf{K}, \mathbf{V} \in \mathbb{R}^{N \times d}$ stored in HBM, a block size parameter for keys $B_c$ where $T_c = \lceil \frac{N}{B_c} \rceil$.
      \STATE Allocate registers according to the designated count, determined by the active consumer warps.
      \STATE Initialize local memory with $\mathbf{O}_i = (0) \in \mathbb{R}^{B_r \times d}$, $\ell_i = (0) \in \mathbb{R}^{B_r}$, and $m_i = (-\infty) \in \mathbb{R}^{B_r}$.
      \STATE Synchronize until both $\mathbf{Q}_i$ and the first block $\mathbf{K}_0$ are present in shared memory.
      \STATE Apply WGMMA to compute $\mathbf{S}_{\mathrm{cur}} = \mathbf{Q}_i \mathbf{K}_0^T$, then commit and ensure completion.
      \STATE Free buffer stage 0 associated with $\mathbf{K}$.
      \STATE Update $m_i$, $\tilde{\mathbf{P}}_{\mathrm{cur}}$, and $\ell_i$ according to $\mathbf{S}_{\mathrm{cur}}$, and adjust the scale of $\mathbf{O}_i$ appropriately.
      \FOR{$1 \le j < T_c - 1$}
        \STATE Wait until $\mathbf{K}_j$ becomes available in shared memory.
        \label{code:mainloop_start}
        \STATE Use WGMMA to calculate $\mathbf{S}_{\mathrm{next}} = \mathbf{Q}_i \mathbf{K}_{j}^T$, commit but allow processing to continue. \label{code:first_wgmma}
        \STATE Hold until $\mathbf{V}_{j-1}$ is present in shared memory.
        \STATE Accumulate $\mathbf{O}_{i}$ with $\tilde{\mathbf{P}}_{\mathrm{cur}} \mathbf{V}_{j-1}$ using WGMMA, committing but not blocking execution. \label{code:second_wgmma}
        \STATE Wait for the completion of the WGMMA calculation $\mathbf{Q}_i \mathbf{K}_{j}^T$.
        \STATE Calculate updated $m_{i}$, $\tilde{\mathbf{P}}_{\mathrm{next}}$, and $\ell_{i}$ based on $\mathbf{S}_{\mathrm{next}}$. \label{code:softmax}
        \STATE Once WGMMA $\tilde{\mathbf{P}}_{\mathrm{cur}} \mathbf{V}_{j-1}$ completes, rescale $\mathbf{O}_i$ accordingly.
        \STATE Release buffer stage $(j\,\%\,s)$ for $\mathbf{K}$ and $(j-1\,\%\,s)$ for $\mathbf{V}$ respectively.
        \STATE Overwrite $\mathbf{S}_{\mathrm{cur}}$ with $\mathbf{S}_{\mathrm{next}}$.
        \label{code:mainloop_end}
      \ENDFOR \STATE Wait until $\mathbf{V}_{T_c - 1}$ is ready in shared memory.
      \STATE Perform
      $\mathbf{O}_{i} = \mathbf{O}_{i} + \tilde{\mathbf{P}}_{\mathrm{last}} \mathbf{V}_{T_c - 1}$ using WGMMA, and ensure commitment and completion.

\STATE Finalization: Adjust the scale of $\mathbf{O}_{i}$ in accordance with $m_{i}$. Determine $L_{i}$ using $m_{i}$ and $\ell_{i}$.  
Store $\mathbf{O}_{i}$ and $L_{i}$ into HBM, assigning them as the $i$-th segments in $\mathbf{O}$ and $L$ respectively.  
\end{algorithmic}
\end{algorithm}

\cref{alg:flash3_wgmma} serves as a substitute for the consumer segment from \cref{alg:flash3_wgmma_ws_only}, thereby assembling the full \textsc{FlashAttention-3} procedure when operating with FP16 data. Conceptually, WGMMA denotes asynchronous GEMM throughout this context. In the core loop, spanning from lines \ref{code:mainloop_start} to \ref{code:mainloop_end}, the execution of the second WGMMA operation for iteration $j$ (line \ref{code:second_wgmma}) is performed concurrently with the softmax computation associated with iteration $j+1$ (line \ref{code:softmax}).

Although the pipelined approach depicted above promises substantial theoretical speedups, several practical factors influence its effectiveness:

\paragraph{Compiler reordering} While the provided pseudocode models an ideal order of operations, the actual instruction sequence may diverge due to compiler (NVCC) optimization. This rearrangement can interfere with the intended interleaving of WGMMA and non-WGMMA instructions, possibly undermining the anticipated performance and pipeline behavior. Examination of the generated SASS code, however, confirms that expected overlapping does occur (see Section~\ref{sec:2-stage-sass}).

\paragraph{Register pressure} To achieve maximum efficiency, avoiding register spilling is crucial. The 2-stage pipeline requires storing more intermediate data in registers to sustain computation across stages, notably an additional $\mathbf{S}_{\mathrm{next}}$ that occupies $B_r \times B_c \times \text{sizeof}(\text{float})$ of register space per threadblock. This increases register consumption and may compete with optimizations that use large block sizes, both of which demand high register availability. Careful profiling is necessary to strike a balance between these competing resource requirements.

\paragraph{3-stage pipelining} Building upon the prior 2-stage method, a 3-stage version can further overlap the second WGMMA with softmax, potentially raising the utilization of Tensor Cores. However, adding an extra pipeline stage heightens register requirements, intensifying the trade-off between pipeline depth and tile size. Further details and empirical evaluation of the 3-stage pipeline are presented in ~\cref{sec:3-stage}.

\subsection{Low-precision with FP8}
\label{sec:algofp8}

\textbf{Efficiency: layout transformations.} Executing the forward pass of \textsc{FlashAttention-3} using FP8 precision introduces unique obstacles regarding layout compatibility that are not present when operating with FP16.

To begin, it is important to observe that the input tensors $\mathbf{Q}$, $\mathbf{K}$, and $\mathbf{V}$ are conventionally arranged contiguously along the head dimension. However, to comply with the k-major requirement of FP8 WGMMA in the second GEMM, $\mathbf{V}$—specifically, the SMEM-loaded tiles of $\mathbf{V}$—must be laid out contiguously along the sequence length axis. Because the TMA load operation cannot modify the contiguous dimension, this necessitates either (1) performing a transpose of $\mathbf{V}$ within GMEM before computation begins, or (2) conducting an in-kernel transpose for each tile of $\mathbf{V}$ once loaded into SMEM. For the first approach, possible implementations include: (1a) integrating the transpose step with the epilogue of a prior operation (such as the rotary embedding), or (1b) invoking a dedicated transpose kernel as a preprocessing routine\footnote{A high-performance transpose kernel can operate at speeds approaching the device's bandwidth~\citep{colfax_cutlass_transpose_2024}.} to swap the strides associated with the sequence length and head dimensions. Unfortunately, (1a) presents significant challenges for inclusion in standard libraries, while (1b) is highly inefficient in situations where inference is constrained by memory bandwidth.

For FP8 \textsc{FlashAttention-3}, we select approach (2). Within the kernel, transposition is implemented via LDSM (\verb|ldmatrix|) and STSM (\verb|stmatrix|) instructions, which coordinate among a warp of threads to transfer data between SMEM and RMEM in 128-byte chunks.\footnote{LDSM/STSM operations are outlined in the PTX documentation as handling $8 \times 8$ blocks consisting of 16-bit values \cite[\S 9.7.13.4.15-16]{ptx}. For FP8, we pack pairs of 8-bit elements together to utilize these instructions. However, the LDSM/STSM transpose variants are unable to separate these paired 8-bit items, which requires additional register operations between the LDSM and STSM steps to effect a correct tile-wise transposition. For brevity, these specifics are not discussed.} Employing LDSM/STSM allows efficient use of registers—so we use them within the producer warpgroup—and also supports format transposition during memory operations. Additionally, once the initial iteration is complete, the next tile-wise transpose of $\mathbf{V}$ can be pipelined to execute concurrently with the two WGMMAs acting on the previous $\mathbf{V}$ tile and the current $\mathbf{K}$ tile, utilizing parallel computation effectively.

Secondly, we note that the FP32 accumulator’s memory layout in an FP8 WGMMA diverges from the register arrangement expected for its operand A, contrasting with the FP16 case. Illustrative segments of these layouts are shown in \cref{fig:rmem_accum} and \cref{fig:rmem_operand}, which map how threads hold register entries in the designated order. Utilizing byte permute instructions, we are able to reformat the accumulator output from the initial WGMMA to be compatible for input to the subsequent WGMMA, ensuring alignment with the $\mathbf{V}$ tile organization that arises from the in-kernel transpose. More concretely, referring to \cref{fig:rmem_accum}, we reorder the register sequence as follows:
$$\{ \verb|d0 d1 d4 d5 d2 d3 d6 d7| \},$$
and this register pattern is repeated for each block of 8 bytes. With respect to the logical structure of the $\mathbf{P}$ tile, this operation effectively permutes the columns (for instance, columns $0189$ are transformed into the leading four columns). To guarantee correct tile computation in WGMMA, the in-kernel transpose can be made to output the $\mathbf{V}$ tile with a corresponding row permutation.\footnote{By permitting such flexibility in the in-kernel transpose, shuffle instructions that reassign register ownership between threads—as previously discussed in~\citep{colfax_fp8_flashattention_2024}—are rendered unnecessary.}

\textbf{Accuracy: block quantization and incoherent processing.} The FP8 (e4m3) representation allocates 3 bits for the mantissa and 4 bits for the exponent, leading to greater numerical error compared with FP16 or BF16. Additionally, when dealing with large-scale models, there are often outlier values~\citep{dettmers2208llm, sun2024massive} with magnitudes far exceeding the majority, complicating the quantization process. A standard solution involves per-tensor scaling~\citep{micikevicius2022fp8}, in which a single scaling factor is assigned to each tensor (such as one for $\mathbf{Q}$, one for $\mathbf{K}$, and one for $\mathbf{V}$). To mitigate numerical inaccuracies in FP8-based attention computations, we introduce the following two methods:

\begin{enumerate}

\item \textbf{Block quantization}: In this approach, a single scalar value is maintained per block; specifically, the tensors corresponding to $\mathbf{Q}$, $\mathbf{K}$, and $\mathbf{V}$ are partitioned into segments of size $B_r \times d$ or $B_c \times d$ and each block is quantized independently. This quantization process can be seamlessly integrated with operations performed immediately before the attention calculation (such as rotary embedding) without introducing additional latency, since rotary embedding is limited by memory bandwidth. Furthermore, as the \textsc{FlashAttention-3} algorithm inherently processes data in block-wise manner, each segment of $\mathbf{S}$ can be rescaled to accommodate block quantization without incurring extra computational overhead.

\item \textbf{Incoherent processing}: To mitigate the influence of anomalous values, both $\mathbf{Q}$ and $\mathbf{K}$ are pre-multiplied by a random orthogonal matrix $\mathbf{M}$ prior to FP8 quantization. Owing to the orthogonality of $\mathbf{M}$, $\mathbf{M} \mathbf{M}^\top = I$, which guarantees that $(\mathbf{Q} \mathbf{M}) (\mathbf{K} \mathbf{M})^\top = \mathbf{Q} \mathbf{K}^\top$, meaning that the multiplication of both $\mathbf{Q}$ and $\mathbf{K}$ by $\mathbf{M}$ preserves the result of the attention mechanism. This transformation helps distribute outlier values, as each component of $\mathbf{Q} \mathbf{M}$ or $\mathbf{K} \mathbf{M}$ represents a random linear combination of the original elements, thereby alleviating quantization errors. For implementation, we adopt the strategy described by \citet{chee2024quip} and~\citet{tseng2024quip}, selecting $\mathbf{M}$ as the product of random diagonal matrices containing $\pm 1$ entries and a Hadamard matrix. This structure permits matrix multiplication in $O(d \log d)$ time, substantially improving over the usual $O(d^2)$ complexity, and it can be efficiently fused with rotary embedding without any extra computational load.

\end{enumerate}

We demonstrate in \cref{sec:numerical_error} that these two methods are capable of decreasing numerical error by as much as 2.6$\times$.

\section{Empirical Validation}
\label{sec:exp}

To construct and assess the performance and precision of \textsc{FlashAttention-3}, we employ key primitives from CUTLASS~\citep{Thakkar_CUTLASS_2023}, specifically utilizing WGMMA and TMA abstractions.

\begin{itemize}

\item \textbf{Benchmarking attention.} We evaluate the performance of \textsc{FlashAttention-3} by recording its runtime over a range of sequence lengths and contrasting these results with the standard PyTorch implementation, \textsc{FlashAttention-2}, a Triton-based version of \textsc{FlashAttention-2} tailored for H100 GPUs, as well as a cuDNN-optimized vendor variant of \textsc{FlashAttention-2} for H100 hardware. Our findings demonstrate that \textsc{FlashAttention-3} achieves up to twice the speed of \textsc{FlashAttention-2}, and operates 1.5$\times$ faster compared to the Triton variant. Moreover, \textsc{FlashAttention-3} attains throughput of up to 740 TFLOPs/s, reaching 75\% of the peak theoretical TFLOPs/s on H100 GPUs.

\item \textbf{Ablation study.} Through an ablation analysis, we establish that the improvements in our algorithm—specifically warp-specialization and GEMM-softmax pipelining—are responsible for the observed increase in speed for \textsc{FlashAttention-3}.

\item \textbf{Accuracy of FP8 attention.} Our experiments further reveal that applying block quantization together with incoherent processing lowers the numerical error of FP8 \textsc{FlashAttention-3} by a factor of 2.6$\times$.
\end{itemize}

\subsection{Benchmarking Attention}
\label{subsec:benchmark_attn}

We evaluate the performance of various attention mechanisms by recording their runtimes on an H100 80GB SXM5 GPU across a range of configurations (including the presence or absence of a causal mask, and head dimensions of either 64 or 128) using FP16 inputs. The benchmark outcomes, presented in~\cref{fig:benchmark_attn_fwd} and~\cref{fig:benchmark_attn_bwd}, indicate that \textsc{FlashAttention-3} achieves a speedup of approximately 1.5--2.0$\times$ over \textsc{FlashAttention-2} during the forward pass, and improves the backward pass speed by 1.5--1.75$\times$. Relative to conventional attention implementations, \textsc{FlashAttention-3} exhibits a substantial speed advantage, running 3--16$\times$ faster. Notably, for sequences of medium to large length (1k tokens or more), \textsc{FlashAttention-3} outperforms a proprietary vendor library (cuDNN -- closed source) specifically tailored for H100 GPUs.

\paragraph{Benchmark settings:} We vary the sequence length as 512, 1k, ..., 16k, and set batch size so that the total number of tokens is 16k. We set the hidden dimension to 2048, and head dimension to be either 64, 128, or 256 (i.e., 32 heads, 16 heads, or 8 heads). To calculate the FLOPs of the forward pass, we use:

\begin{equation*}
  4 \cdot \text{seqlen}^2 \cdot \text{head dimension} \cdot \text{number of heads}.
\end{equation*}

Applying causal masking entails halving the computed value, as typically just half of the elements are processed. The computation for the backward pass FLOPs is obtained by scaling the forward pass FLOPs by 2.5; this reflects the presence of 2 matrix multiplications during the forward pass versus 5 matrix multiplications required in the backward pass because of recomputation.

Additionally, we evaluate the FP8 runtime during the forward pass using comparable configurations. The findings for a headdim of 256 are presented in ~\cref{fig:benchmark_attn_fp8}, while comprehensive results are detailed in \cref{sec:benchmark_attn_fp8_full}.

\subsection{Ablation Study: 2-Stage Pipelining Experiments}
\label{sec:2-stage-pipelining-experiments}

We conduct an ablation study on both the two-stage WGMMA-softmax pipelining and the use of warp-specialization for the non-causal FP16 \textsc{FlashAttention-3}, holding the parameters $\{\text{batch}, \text{seqlen}, \text{nheads}, \text{hdim}\} = \{ 4, 8448, 16, 128\}$ constant.
As presented in~\cref{table:ablation_pipelining}, these experiments demonstrate that integrating asynchronous execution with warp-specialization and facilitating overlap between GEMM and softmax computations yields a substantial improvement in performance, boosting throughput from 570 to 661 TFLOPs.

\subsection{Numerical Error Validation}
\label{sec:numerical_error}

Due to increasing concerns regarding numerical inaccuracies in \textsc{FlashAttention}~\citep{golden2024flash}, we evaluate \textsc{FlashAttention-2}, \textsc{FlashAttention-3}, and the conventional attention mechanism, benchmarking each against a reference FP64 implementation. In order to mimic extreme feature values and activation patterns seen in LLMs~\citep{dettmers2208llm, sun2024massive}, the matrices $\mathbf{Q}, \mathbf{K}, \mathbf{V}$ are populated according to the subsequent distribution:

\begin{equation*}
  \mathcal{N}(0, 1) + \mathcal{N}(0, 100) \cdot \mathrm{Bernoulli}(0.001).
\end{equation*}

Specifically, each matrix element is sampled from a normal distribution with zero mean and unit standard deviation, except for 0.1\% of the entries, which have an additional independent term drawn from a normal distribution with standard deviation 10. The root mean squared error (RMSE) is reported in~\cref{table:numerical_error}.
In FP16, \textsc{FlashAttention-2} and \textsc{FlashAttention-3} both yield an RMSE that is 1.7$\times$ lower relative to the conventional approach, owing to the use of FP32 for storing intermediate softmax results. The baseline attention in FP8 mode applies per-tensor scaling, uses FP32 for matmul accumulation, and maintains intermediate softmax outputs in FP16. Owing to its use of block quantization and non-sequential processing, \textsc{FlashAttention-3} in FP8 demonstrates a 2.6$\times$ improvement in accuracy over this baseline.

\section{Dicussion, Limitations, Conclusion}
\label{sec:discussion}

Through our development of \textsc{FlashAttention-3}, we illustrate that innovations in programming strategies and hardware functionalities—specifically asynchronous execution and low-precision computation—can substantially enhance both performance and accuracy in attention mechanisms. Our solution accelerates attention operations by 1.5--2.0$\times$ relative to \textsc{FlashAttention-2}, and mitigates FP8 numerical error by a factor of 2.6$\times$ when compared to typical per-tensor quantization schemes. Looking forward, several areas remain for improvement: tailoring the system for large language model (LLM) inference, incorporating a persistent kernel architecture into the FP8 implementation,\footnote{In our benchmarking process, FP16 \textsc{FlashAttention-3} utilizes a persistent kernel and employs load balancing, whereas FP8 \textsc{FlashAttention-3} lacks these features. This, in part, accounts for the inferior performance of FP8 \textsc{FlashAttention-3} at short sequence lengths and with causal masking, relative to FP8 cuDNN kernels.} and further exploring how low-precision attention affects large-scale model training. Although our experiments center on Hopper GPUs, we anticipate that these methodological advancements will be transferable across a broader spectrum of hardware accelerators. Ultimately, we envision that improved primitives like attention will foster the development of novel applications requiring extended contextual understanding.

\end{document}